%!TEX program = xelatex
\def\numeromodulo{2}
\def\numerolista{2}
\documentclass[12pt]{exam}

\def\ano{2023}
\def\semestre{1}
\def\disciplina{Introdução à \'Algebra Linear}
\def\turma{3}

\usepackage{caption}
\usepackage{amssymb}
\usepackage{amsmath,amsfonts,amsthm,amstext}
\usepackage[utf8]{inputenc} % Required for inputting international characters
\usepackage[T1]{fontenc}
\usepackage[brazil]{babel}
\usepackage{graphicx}
\graphicspath{{../figuras_geogebra/}{../Pictures/}}
\usepackage{enumitem}
\usepackage{multicol}
\usepackage{answers}
\usepackage[svgnames]{xcolor}
\usepackage{tikz}
\usepackage{ifthen}
\usetikzlibrary{lindenmayersystems}
\usetikzlibrary[shadings]
\usepackage{tikz-3dplot}
\usepackage{tcolorbox}
%-------------amatrix
% Augmented matrix.  Usage (note the argument does not count the aug col):
% \begin{amatrix}{2}
%   1  2  3 \\  4  5  6
% \end{amatrix}
\newenvironment{amatrix}[1]{%
  \left(\begin{array}{@{}*{#1}{r}|r@{}}
}{%
  \end{array}\right)
}

\Newassociation{solucao}{Solution}{ans}
\newtheorem{exercicio}{}

\setlength{\topmargin}{-1.0in}
\setlength{\oddsidemargin}{0in}
\setlength{\textheight}{10.1in}
\setlength{\textwidth}{6.5in}
\setlength{\baselineskip}{12mm}

\newcounter{exercicios}
\setcounter{exercicios}{0}
\newcommand{\questao}{
\addtocounter{exercicios}{1}
\noindent{\bf Exerc{\'\i}cio \Roman{exercicios}: }}

\newcommand{\resp}[1]{
\noindent{\bf Exerc{\'\i}cio #1: }}

\extraheadheight{0.7in}
\firstpageheadrule
\runningheadrule
\lhead{
        \begin{minipage}[c]{1.7cm}
        \includegraphics[width=1.7cm]{unb.pdf}
        \end{minipage}%
        \hspace{0pt}
        \begin{minipage}[c]{4in}
          {Universidade de Bras{\'\i}lia} --
          {Departamento de Matem{\'a}tica}
\end{minipage}
\vspace*{-0.8cm}
}
% \chead{Universidade de Bras{\'\i}lia - Departamento de Matem{\'a}tica}
% \rhead{}
% \vspace*{-2cm}

\extrafootheight{.5in}
\footrule
\lfoot{\disciplina\ - \semestre$^o$/\ano\ - Módulo \numeromodulo}
\cfoot{}
\rfoot{P\'agina \thepage\ de \numpages}
\newcommand{\nsub}{\varsubsetneq}
\newcommand{\vaz}{\emptyset}
\newcommand{\im}{{\rm Im\,}}
\newcommand{\sub}{\subseteq}
\renewcommand{\sin}{{\rm sen\,}}
\newcommand{\n}{\mathbb{N}}
\newcommand{\z}{\mathbb{Z}}
\newcommand{\rac}{\mathbb{Q}}
\newcommand{\real}{\mathbb{R}}
\newcommand{\complex}{\mathbb{C}}
\newcommand{\cp}[1]{\mathbb{#1}}
\newcommand{\ch}{\mbox{\textrm{car\,}}\nobreak}
\newcommand{\vesp}[1]{\vspace{ #1  cm}}
\newcommand{\compcent}[1]{\vcenter{\hbox{$#1\circ$}}}
\newcommand{\comp}{\mathbin{\mathchoice
{\compcent\scriptstyle}{\compcent\scriptstyle}
{\compcent\scriptscriptstyle}{\compcent\scriptscriptstyle}}}

\newcommand{\integer}{\mathbb{Z}}
\begin{document}

  \Opensolutionfile{ans}[ans1]
  \begin{center}
    {\Large\bf \disciplina\ - Turma \turma\ -- \semestre$^{o}$/\ano} \\ \vspace{9pt} {\large\bf
        $\numerolista^a$ Lista de Exercícios -- Módulo \numeromodulo}\\ \vspace{9pt} Prof. José Antônio O. Freitas
  \end{center}
  \hrule

\vesp{.6}

\begin{exercicio}
    Quais dos seguintes vetores são combinações lineares de $u = (0, -2, 2)$ e $v = (1, 3, -1)$?
    \begin{multicols}{2}
        \begin{enumerate}[label={\alph*})]
            \item $w = (2, 2, 3)$

            \item $w = (3, 1, 5)$

            \item $w = (0, 4, 5)$

            \item $w = (0, 0, 0)$
        \end{enumerate}
    \end{multicols}
    \begin{solucao}
      \begin{multicols}{2}
        \begin{enumerate}[label={\alph*})]
          \item Não.

          \item Sim.

          \item Não.

          \item Sim.
        \end{enumerate}
      \end{multicols}
    \end{solucao}
\end{exercicio}

\begin{exercicio}
    Expresse os seguintes vetores como combinações lineares de $u = (2, 1, 4)$, $v = (1, -1, 3)$ e $w = (3, 2, 5)$.
    \begin{multicols}{2}
        \begin{enumerate}[label={\alph*})]
            \item $(-9, -7, -15)$

            \item $(6, 11, 6)$

            \item $(0, 0, 0)$

            \item $(7, 8, 9)$
        \end{enumerate}
    \end{multicols}
\end{exercicio}

\begin{exercicio}
    Quais dos seguintes vetores são combinações lineares de:
    \[
        A = \begin{bmatrix}
            \phantom{-} 4 & 0\\-2 & 2
        \end{bmatrix},
        B = \begin{bmatrix}
            1 & -1\\2 & \phantom{-} 3
        \end{bmatrix},
        C = \begin{bmatrix}
            0 & 2\\1 & 4
        \end{bmatrix}.
    \]
    \begin{multicols}{2}
        \begin{enumerate}[label={\alph*})]
            \item $\begin{bmatrix}
                \phantom{-} 6 & -8\\-1 & -8
            \end{bmatrix}$

            \item $\begin{bmatrix}
                0 & 0\\0 & 0
            \end{bmatrix}$

            \item $\begin{bmatrix}
                6 & 0\\3 & 8
            \end{bmatrix}$

            \item $\begin{bmatrix}
                -1 & 5\\\phantom{-} 7 & 1
            \end{bmatrix}$
        \end{enumerate}
    \end{multicols}
    \begin{solucao}
      \begin{multicols}{2}
        \begin{enumerate}[label={\alph*})]
          \item Não.

          \item Sim.

          \item Não

          \item Não.
        \end{enumerate}
      \end{multicols}
    \end{solucao}
\end{exercicio}

\begin{exercicio}
    Em cada caso, expresse o vetor dado como uma combinação linear de
    \[
        p_1 = 2 + x + 4x^2,
        p_2 = 1 - x + 3x^2,
        p_3 = 3 + 2x + 5x^2.
    \]
    \begin{multicols}{2}
        \begin{enumerate}[label={\alph*})]
            \item $q = -9 -7x - 15x^2$

            \item $q = 6 + 11x + 6x^2$

            \item $q = 0$

            \item $q = 7 + 8x + 9x^2$
        \end{enumerate}
    \end{multicols}
\end{exercicio}

\begin{exercicio}
    Sejam $v_1 = (2, 1, 0, 3)$, $v_2 = (3, -1, 5, 2)$ e $v_3 = (-1, 0, 2, 1)$. Em cada caso decida se o vetor dado pertence à $W = [v_1, v_2, v_3]$.
    \begin{multicols}{2}
        \begin{enumerate}[label={\alph*})]
            \item $(2, 3, -7, 3)$

            \item $(0, 0, 0, 0)$

            \item $(1, 1, 1, 1)$

            \item $(-4, 6, -13, 4)$

            \item $(4, 6, -13, 4)$

            \item $(1, 0, 1, 1)$
        \end{enumerate}
    \end{multicols}
    \begin{solucao}
      \begin{multicols}{3}
        \begin{enumerate}[label={\alph*})]
          \item Sim.

          \item Sim.

          \item Não.

          \item Sim.

          \item Não.

          \item Sim.
        \end{enumerate}
      \end{multicols}
    \end{solucao}
\end{exercicio}

Nos exercícios \ref{subespacoinicio} à \ref{subespacofim}, verifique se $S$ é um subespaço vetorial do espaço vetorial $V$ sobre o corpo $\cp{K}$ em questão:
\begin{exercicio}\label{subespacoinicio}
  $V = \real^n$ e $S = \{(a_1, a_2, \dots, a_n) \in \real^n \mid a_1a_2 = 0\}$; $\cp{K} = \real$.
  \begin{solucao}
    Não é subespaço.
  \end{solucao}
\end{exercicio}

\begin{exercicio}
    $V = \real^3$ e $S = \left\{\begin{bmatrix}a\\0\\a\end{bmatrix} \mid a \in \real\right\}$; $\cp{K} = \real$.
    \begin{solucao}
      É subespaço.
    \end{solucao}
\end{exercicio}

\begin{exercicio}
    $V = \real^3$ e $S = \left\{\begin{bmatrix}a\\-a\\2a\end{bmatrix} \mid a \in \real\right\}$; $\cp{K} = \real$.
    \begin{solucao}
      É subespaço.
    \end{solucao}
\end{exercicio}

\begin{exercicio}
    $V = \real^3$ e $S = \left\{\begin{bmatrix}a\\b\\a + b + 1\end{bmatrix} \mid a \in \real\right\}$; $\cp{K} = \real$.
    \begin{solucao}
      Não é subespaço.
    \end{solucao}
\end{exercicio}

\begin{exercicio}
    $V = \mathcal{P}_2(\rac)$ e $S = \{a + bx + cx^2 \mid a + b + c = 0\}$; $\cp{K} = \rac$.
    \begin{solucao}
      É subespaço.
    \end{solucao}
\end{exercicio}

\begin{exercicio}
    $V = \mathcal{P}(\real)$ e $S = \{a + bx + cx^2 + dx^3\mid a, b, c, d \in \real\}$; $\cp{K} = \rac$.
    \begin{solucao}
      É subespaço.
    \end{solucao}
\end{exercicio}

\begin{exercicio}
  $V = \cp{M}_2(\complex)$ e $S = \left\{\begin{pmatrix} a_{11} & a_{12}\\ a_{21} & a_{22}\end{pmatrix} \in V \mid a_{ij} = \overline{a_{ji}}, i, j = 1, 2\right\}$; $\cp{K} = \complex$, onde $\overline{a + bi} = a - bi$.
  \begin{solucao}
    Não é subespaço.
  \end{solucao}
\end{exercicio}

\begin{exercicio}
  $V = \real^3$ e $S = \{(x_1, x_2, x_1x_2) \in \real^3 \mid x_1, x_2 \in \real\}$; $\cp{K} = \real$.
  \begin{solucao}
    Não é subespaço.
  \end{solucao}
\end{exercicio}

\begin{exercicio}
  $V = \real^2$ e $S = \{(x, y) \in \real^2 \mid x + 3y = 0\}$; $\cp{K} = \real$.
  \begin{solucao}
    é subespaço.
  \end{solucao}
\end{exercicio}

\begin{exercicio}
  $V = \real^2$ e $S = \{(x, y) \in \real^2 \mid y = x + 1\}$; $\cp{K} = \real$.
  \begin{solucao}
    Não é subespaço.
  \end{solucao}
\end{exercicio}

\begin{exercicio}
  $V = \real^2$ e $S = \{(x, y) \in \real^2 \mid x \ge 0\}$; $\cp{K} = \real$.
  \begin{solucao}
    Não é subespaço.
  \end{solucao}
\end{exercicio}

\begin{exercicio}
  $V = \real^3$ e $S = \{(x, y, z) \in \real^3 \mid y = x + 2 \mbox{ e } z = 0\}$; $\cp{K} = \real$.
  \begin{solucao}
    Não é subespaço.
  \end{solucao}
\end{exercicio}

\begin{exercicio}
  $V = \real^3$ e $S = \{(x, y, z) \in \real^3 \mid x + 2y -3z = 4\}$; $\cp{K} = \real$.
  \begin{solucao}
    Não é subespaço.
  \end{solucao}
\end{exercicio}

\begin{exercicio}
  $V = \real^3$ e $S = \{(x, y, z) \in \real^3 \mid x = -z \mbox{ e } x = z\}$; $\cp{K} = \real$.
  \begin{solucao}
    é subespaço.
  \end{solucao}
\end{exercicio}

\begin{exercicio}
  $V = \real^3$ e $S = \left\{(x, y, z) \in \real^3 \mid \dfrac{x}{2} = \dfrac{y - 3}{5}\right\}$; $\cp{K} = \real$.
  \begin{solucao}
    Não é subespaço.
  \end{solucao}
\end{exercicio}

\begin{exercicio}
  $V = \real^4$ e $S = \{(x_1, x_2, x_3, x_4) \in \real^4 \mid x_1 + x_2 = 0, x_3 - x_4 = 0 \in \real\}$; $\cp{K} = \real$.
  \begin{solucao}
    é subespaço.
  \end{solucao}
\end{exercicio}

\begin{exercicio}
  $V = \real^4$ e $S = \{(x_1, x_2, x_3, x_4) \in \real^4 \mid 2x_1 + x_2 - x_4 = 0, x_3 = 0 \in \real\}$; $\cp{K} = \real$.
  \begin{solucao}
    é subespaço.
  \end{solucao}
\end{exercicio}

\begin{exercicio}
  $V = \cp{M}_2(\complex)$ e $S = \left\{\begin{bmatrix} a & b\\ c & d\end{bmatrix} \in V \mid b = c\right\}$; $\cp{K} = \complex$.
  \begin{solucao}
    é subespaço.
  \end{solucao}
\end{exercicio}

\begin{exercicio}
  $V = \cp{M}_2(\real)$ e $S = \left\{\begin{bmatrix} a & b\\ c & d\end{bmatrix} \in V \mid c = a + b \mbox{ e } d = 0\right\}$; $\cp{K} = \real$.
  \begin{solucao}
    é subespaço.
  \end{solucao}
\end{exercicio}

\begin{exercicio}
  $V = \cp{M}_2(\real)$ e $S = \left\{\begin{bmatrix} a & b\\ 0 & c\end{bmatrix} \in V \mid a, b, c \in \real\right\}$; $\cp{K} = \real$.
  \begin{solucao}
    é subespaço.
  \end{solucao}
\end{exercicio}

\begin{exercicio}
  $V = \cp{M}_2(\real)$ e $S = \left\{\begin{bmatrix} a & b\\ b & c\end{bmatrix} \in V \mid a, b, c \in \real\right\}$; $\cp{K} = \real$.
  \begin{solucao}
    é subespaço.
  \end{solucao}
\end{exercicio}

\begin{exercicio}\label{subespacofim}
  $V = \cp{M}_2(\real)$ e $S = \left\{\begin{bmatrix} a & a+b\\ a-b & b\end{bmatrix} \in V \mid a, b \in \real \right\}$; $\cp{K} = \real$.
  \begin{solucao}
    é subespaço.
  \end{solucao}
\end{exercicio}

\begin{exercicio}
  Sejam $W_1$ e $W_2$  subespaços de um $\cp{K}$-espaço vetorial $V$. Defina
  \[
  W_1 - W_2 = \{ u_1 - u_2 \mid u_1 \in W_1, u_2 \in W_2\}.
  \]
  O conjunto $W_1 - W_2$ é um subespaço vetorial de $V$?
\end{exercicio}

\begin{exercicio}
    Sejam $U$ e $V$ espaços vetoriais sobre um mesmo corpo $\cp{K}$. Defina o \textbf{produto cartesiano} de $U$ e $V$ como sendo
    \[
        U \times V = \{(x, y) \mid x \in U, y \in V\}.
    \]
    Prove que $U \times V$ é um espaço vetorial sobre $\cp{K}$.
\end{exercicio}

\begin{exercicio}
    Seja $W$ um subespaço vetorial de um $\cp{K}$-espaço vetorial $V$. Prove que $\Delta = \{(x, x) \mid x \in W\}$ é um subespaço de $V \times V$.
\end{exercicio}

\begin{exercicio}
  Seja $W = \left\{\begin{pmatrix} a_{11} & a_{12}\\ a_{21} & a_{22}\end{pmatrix} \in \cp{M}_2(\complex) \mid a_{11} + a_{12} = 0\right\}$.
  \begin{enumerate}[label={\alph*})]
    \item Mostre que $W$ é um espaço vetorial sobre $\real$.

    \item Seja $W_1 = \{(a_{ij})_{i,j}, \in \cp{M}_2(\complex), i, j = 1, 2 \mid a_{21} = -\overline{a_{12}}\}$, onde $\overline{a + bi} = a - bi$. Prove que $W_1$ é um subespaço de $\cp{M}_2(\complex)$ sobre $\real$.
  \end{enumerate}
\end{exercicio}

\begin{exercicio}
  Considere o subespaço $S = [(1,1,-2,4),(1,1,-1,2),(1,4,-4,8)]$ de $\real^4$.
  \begin{enumerate}[label={\alph*})]
    \item O vetor $(2/3, 1, -1, 2)$ pertence a $S$?

    \item O vetor $(0, 0, 1, 1)$ pertence a $S$?

    \item O vetor $(1, 0, 0, 0)$ pertence a $S$?

    \item O vetor $(0, 1, 0, 0)$ pertence a $S$?

    \item O vetor $(0, 0, 1, 0)$ pertence a $S$?

    \item O vetor $(0, 0, 0, 1)$ pertence a $S$?
  \end{enumerate}
  \begin{solucao}
    \begin{multicols}{3}
      \begin{enumerate}[label={\alph*})]
        \item Sim.

        \item Não.

        \item Sim.

        \item Sim.

        \item Não.

        \item Não.
      \end{enumerate}
    \end{multicols}
  \end{solucao}
\end{exercicio}

\begin{exercicio}
    Considere o $\rac$-espaço vetorial $V = M_2(\rac)$. Sejam
    \[
        A = \begin{bmatrix}
            \phantom{-} 1 & 1\\
            -1 & 1
        \end{bmatrix},
        B = \begin{bmatrix}
            1 & -1\\
            1 & \phantom{-} 0
        \end{bmatrix}
    \]
    e $W = Span(A, B)$. Determine se a matriz $C$ está em $W$ nos seguintes casos:
    \begin{multicols}{2}
      \begin{enumerate}[label={\alph*})]
          \item $C = \begin{bmatrix}1 & 2\\3 & 4\end{bmatrix}$

          \item $C = \begin{bmatrix}3 & -5\\5 & \phantom{-} 1\end{bmatrix}$

          \item $C = \begin{bmatrix}-1 & -5\\ \phantom{-} 5 & -1\end{bmatrix}$

          \item $C = \begin{bmatrix}0 & -1\\1 & \phantom{-} 0\end{bmatrix}$
      \end{enumerate}
    \end{multicols}
    \begin{solucao}
      \begin{multicols}{2}
        \begin{enumerate}[label={\alph*})]
          \item Não.

          \item Sim.

          \item Não.

          \item Sim.
        \end{enumerate}
      \end{multicols}
    \end{solucao}
\end{exercicio}

\begin{exercicio}
  Seja $W$ o subespaço de $M_2(\real)$ definido por
  \[
  W = \left\{\begin{bmatrix}2a & a + 2b\\0 & a - b\end{bmatrix} \mid a, b \in \real\right\}.
  \]
  \begin{enumerate}[label={\alph*})]
    \item $\begin{bmatrix}0 & -2\\0 & \phantom{-} 1\end{bmatrix} \in W$?

    \item $\begin{bmatrix}0 & 2\\3 & 1\end{bmatrix} \in W$?
  \end{enumerate}
  \begin{solucao}
    \begin{enumerate}[label={\alph*})]
      \item Sim.

      \item Não.
    \end{enumerate}
  \end{solucao}
\end{exercicio}

\begin{exercicio}
  Considere o subespaço de $\real^4$ gerado pelos vetores $v_1 = (1, -1, 0, 0)$, $v_2 = (0, 0, 1, 1)$, $v_3 = (2, -2, 1, 1)$ e $v_4 = (1, 0, 0, 0)$.
  \begin{enumerate}[label={\alph*})]
    \item O vetor $(2, -3, 2, 2) \in [v_1, v_2, v_3, v_4]$? Justifique.

    \item O vetor $(1, 2, 3, 4) \in [v_1, v_2, v_3, v_4]$? Justifique.

    \item O vetor $(0, -3, 1, -2) \in [v_1, v_2, v_3, v_4]$? Justifique.

    \item O vetor $(1, 0, 0, 1) \in [v_1, v_2, v_3, v_4]$? Justifique.

    \item O vetor $(0, 1, 0, 0) \in [v_1, v_2, v_3, v_4]$? Justifique.

    \item O vetor $(0, 0, 1, 0) \in [v_1, v_2, v_3, v_4]$? Justifique.

    \item O vetor $(0, 0, 0, 1) \in [v_1, v_2, v_3, v_4]$? Justifique.
  \end{enumerate}
  \begin{solucao}
    \begin{multicols}{3}
      \begin{enumerate}[label={\alph*})]
        \item Sim.

        \item Sim.

        \item Sim.

        \item Sim.

        \item Sim.

        \item Sim.

        \item Sim.
      \end{enumerate}
    \end{multicols}
  \end{solucao}
\end{exercicio}

\newpage
\Closesolutionfile{ans}

\begin{center}
{\large\bf RESPOSTAS}
\end{center}
\hrule
\input{ans1}

\end{document}
